\section{TEMARIO, OBJETIVOS Y PROGRAMA}

\subsection{MOTIVACIÓN}
\textbf{Duración:} 1 hr\\
Motivar al grupo respecto a la importancia de la protección radiológica y de la necesidad de que la seguridad constituya un hábito personal. Explicará la responsabilidad adquirida al manejar una fuente radioactiva.

\subsection{ESTRUCTURA ATÓMICA: Terminología}
\textbf{Duración:} 1 hr\\
El alumno identificará definiciones de: átomo, núcleo, protón, electrón, número atómico, número de masa, número de neutrones; explicará correctamente estas definiciones y la simbología $A^Z$.

\subsection{RADIOACTIVIDAD Y RADIACIÓN}
\textbf{Duración:} 1:30 hr\\
El alumno podrá describir el fenómeno de la desintegración radioactiva y las características de las radiaciones alfa, beta y gamma. Explicará el concepto de vida media y podrá estimar la actividad inicial y en un tiempo considerado de una media, activada actual de una fuente complemento, empleando una gráfica o una tabla de decaimiento. Empleará correctamente las unidades de actividad: Curie y Becquerel.

\subsection{INTERACCIÓN DE LA RADIACIÓN CON LA MATERIA}
\textbf{Duración:} 1 hr\\
Definirá con sus palabras cada uno de los siguientes conceptos: ion, ionización, excitación, absorción, alcance o penetración, dispersión, radiación primaria, radiación dispersa y radiación transmitida.

\subsubsection{PRÁCTICA NO. 1.}
\textbf{Duración:} 1 hr\\
Absorción de las radiaciones beta y gamma.

\subsection{EXPOSICIÓN Y DOSIS. UNIDADES}
\textbf{Duración:} 1 hr\\
El alumno explicará la diferencia entre exposición y dosis, exposición y rapidez de exposición, dosis y rapidez de dosis. Definirá Roentgen, el rad y sus submúltiplos, así como las unidades de rapidez de exposición y de dosis. Describirá los conceptos de equivalente de dosis y rem, así como las unidades Gray y Sievert y sus equivalentes a rad y rem.